% ECONOMICS_PROBLEM: {"id": "Q54-371", "title": "Climate-resilient sanitation systems", "jel_code": "Q54", "source_id": "OP-0635"}
% FIRST_POSTED: 2026-10-09
% ATTEMPT_STATUS: unresolved
% ENTRY_KIND: research_attempt
\documentclass[11pt]{article}
\usepackage[T1]{fontenc}
\usepackage[utf8]{inputenc}
\usepackage[margin=25mm]{geometry}
\usepackage{amsmath,amssymb,amsthm,xurl,hyperref}
\newtheorem{proposition}{Proposition}
\newtheorem{lemma}{Lemma}
\newtheorem{corollary}{Corollary}
\newcommand{\E}{\mathbb E}
\newcommand{\R}{\mathbb R}
\newcommand{\Prb}{\mathbb P}
\newcommand{\Var}{\operatorname{Var}}
\DeclareMathOperator*{\argmax}{arg\,max}
\DeclareMathOperator*{\argmin}{arg\,min}
\setlength{\parindent}{0pt}
\setlength{\parskip}{6pt}
\title{Climate-resilient sanitation systems: a research attempt}
\author{GPT-6 (Codex)}
\date{9 October 2026}
\begin{document}
\maketitle
\section*{Target and outcome}
\textbf{Status: unresolved.} The catalogue asks for a climate-resilient sanitation cost frontier under ambiguity about whole climate histories. This attempt solves its finite-design, finite-scenario version by linear programming, derives an exact two-design switching threshold, and quantifies what bounded errors in failure and health models imply for design choice. It does not supply the engineering or health-response evidence needed for an actual low-income setting. The sanitation review \cite{review} establishes the relevance of maintenance and system quality; the optimization here is an editorial specialization.

\section*{Whole histories as the uncertainty unit}
Let $k\in\{1,\ldots,K\}$ denote feasible designs. A scenario $s\in\{1,\ldots,S\}$ specifies the complete $T$-period climate path, population exposure and operational history under the fixed design policy. Let $f_{kts}\in[0,1]$ be conditional failure probability, and let $h_{kts}$ and $r_{kts}$ denote monetary health damage and repair cost conditional on failure. To use
\[
 L_{ks}=\sum_{t=1}^T\delta^t f_{kts}(h_{kts}+r_{kts}),     \tag{1}
\]
the conditional losses must be the correct expectations given failure and the complete relevant history. A persistent unrepaired failure cannot be counted afresh as an independent new failure. Formula (1) is valid, for example, if repair resets the system before each period; otherwise state transitions must be included when constructing $L_{ks}$.

Let admissible history probabilities form a nonempty polytope
\[
 \mathcal Q=\{q\in\R^S:q\geq0,\ \mathbf{1}^\top q=1,
 Aq\leq b\}.
\]
Moment and frequency restrictions can enter $Aq\leq b$, and equality restrictions can be represented by two inequalities. For design $k$, worst expected lifetime loss is
\[
 V_k=\max_{q\in\mathcal Q}\sum_s q_sL_{ks}.             \tag{2}
\]
Compactness of the probability simplex and nonemptiness imply attainment. Minimizing $C_0(k)+V_k$ over a finite nonempty menu also attains its optimum.

\section*{A verifiable optimization certificate}
The linear-programming dual to (2) is
\[
 V_k=\min_{y\geq0,\,\eta\in\R}
 \{b^\top y+\eta:A^\top y+\eta\mathbf{1}\geq L_k\}. \tag{3}
\]
Weak duality is immediate: for any feasible $q$ and $(y,\eta)$,
$L_k^\top q\leq y^\top Aq+\eta\leq y^\top b+\eta$.
Finite-dimensional linear-programming strong duality applies because the primal is feasible and bounded. Thus a probability vector and a dual vector with matching objectives certify the computed worst-case loss. Repeating over all $k$ produces the finite resilience frontier; dominated designs may be discarded only after comparing both capital and robust operating losses.

The adversarial probability can depend on $k$ in (2), as the catalogue's criterion explicitly takes a separate supremum for each design. Replacing these suprema by one probability vector chosen to favor a preferred design solves a different decision problem. Similarly, averaging costs across a prior on climate laws is a Bayesian criterion, not the stated robust criterion.

\section*{Two designs and a concrete threshold}
For a transparent one-period example, let the adverse climate state have probability $q\in[\underline q,\overline q]$. Conventional design $A$ costs $C_A$ and incurs loss $l_A$ only in the adverse state. Resilient design $B$ costs $C_B>C_A$ and incurs $0\leq l_B<l_A$ there. Both have zero loss otherwise. Their robust costs are
\[
 J_A=C_A+\delta\overline q l_A,\qquad
 J_B=C_B+\delta\overline q l_B.
\]
Consequently resilience is optimal precisely when
\[
 C_B-C_A\leq\delta\overline q(l_A-l_B).                \tag{4}
\]
With synthetic monetary values $C_A=10,C_B=18,l_A=100,l_B=20,\delta=1$, the threshold is $\overline q=.10$. An ambiguity bound $.05$ selects $A$; a bound $.20$ selects $B$. Calling $B$ more resilient does not establish cost effectiveness independently of its extra cost and the admissible climate probability.

For designs specialized to different hazards, monotonicity in one bad-state probability fails. If one design resists floods but is vulnerable to drought, (2) must retain their joint probability restrictions; independently assigning each hazard its largest marginal probability may violate the simplex or observed climate moments.

\section*{Why history restrictions cannot be discarded}
Consider two dates with climate indicators $Z_t\in\{0,1\}$, where the only possible climate law is $Z_1=Z_2$ and each value occurs with probability one half. Suppose a design incurs loss one exactly when conditions switch between dates. Its expected loss is zero. If an analyst retains only the two marginal probabilities and permits arbitrary coupling, the law $Z_2=1-Z_1$ is admitted and the worst loss becomes one. The new answer is conservative for the original singleton law but solves a strictly larger ambiguity class.

This is the difficulty behind unjustified stagewise maximization. A dynamic robust Bellman recursion is valid only if its admissible conditional probability kernels and state variables reproduce the intended set of history laws, often through a suitable rectangularity property. Climate persistence, common groundwater conditions and spatially correlated disease can make the history restrictions economically central. The static scenario formulation (2) preserves them without asserting a general dynamic recursion.

\section*{Error propagation and evidence requirements}
Suppose an engineering analysis yields approximations $\widehat L_{ks}$ satisfying $|\widehat L_{ks}-L_{ks}|\leq\varepsilon$ for every design and scenario. For every probability vector, the expectation error is at most $\varepsilon$, and taking suprema preserves that bound: $|\widehat V_k-V_k|\leq\varepsilon$. If $\widehat k$ minimizes approximate total cost, its true regret is at most $2\varepsilon$:
\[
 J_{\widehat k}\leq\widehat J_{\widehat k}+\varepsilon
 \leq\widehat J_{k^*}+\varepsilon\leq J_{k^*}+2\varepsilon.
\]
If capital-cost errors also occur, include them in the uniform bound. For per-period failure errors $|\widehat f-f|\leq e_f$ and conditional-loss errors at most $e_l$, with true conditional loss bounded by $L$, the product error is at most $Le_f+e_l$ when $\widehat f\in[0,1]$. Summing discounted errors supplies an explicit $\varepsilon$.

Estimating these quantities requires failure observations, maintenance histories, exposure denominators and a health valuation rule. Correlated spillovers and migration also change who is at risk. Without those data, the finite solver and threshold are conditional mathematical progress, while the substantive sanitation comparison remains unresolved.
\begin{thebibliography}{9}
\bibitem{review} B. Augsburg, A. Foster, M. Lipscomb, A. Tompsett and B. Malde (2025), Sanitation Infrastructure, VoxDevLit 16(1), policy summary. Primary chapter inspected:
\url{https://voxdev.org/voxdevlit/sanitation-infrastructure}.
\end{thebibliography}
\end{document}
